\documentclass[preview]{standalone}

\usepackage{amsmath}
\usepackage{amssymb}
\usepackage{stellar}
\usepackage{definitions}
\usepackage{bettelini}

\begin{document}

\id{linearalgebra-inner-products}
\genpage

\section{Inner products}

\begin{snippetdefinition}{inner-product-definition}{Inner product}
    Let \(V\) be a \vectorspace over a \field
    \(\mathbb{F}\in\{\realnumbers,\complexnumbers\}\). An
    \emph{inner product} on \(V\) is a \function
    \[
        \langle\cdot,\cdot\rangle\colon V\cartesianprod V\fromto\mathbb{F}
    \]
    such that, for every \(u,v,w\in V\) and every
    \(\lambda\in\mathbb{F}\),
    \begin{enumerate}
        \item \(\langle v,v\rangle\geq0\);
        \item \(\langle v,v\rangle=0\) \ifandonlyif \(v=0\);
        \item \(\langle u+v,w\rangle=\langle u,w\rangle+\langle v,w\rangle\);
        \item \(\langle \lambda u,v\rangle=\lambda\langle u,v\rangle\);
        \item \(\langle u,v\rangle=\overline{\langle v,u\rangle}\).
    \end{enumerate}
    If \(\mathbb{F}=\realnumbers\), the last condition becomes
    \[
        \langle u,v\rangle=\langle v,u\rangle.
    \]
\end{snippetdefinition}

\begin{snippetproposition}{real-inner-product-second-variable-linearity}{Linearity in the second variable over the real numbers}
    Let \(V\) be a \vectorspace over \(\realnumbers\), and let
    \(\langle\cdot,\cdot\rangle\) be an \innerproduct on \(V\). Then, for
    every \(u,v,w\in V\) and every \(\lambda\in\realnumbers\),
    \[
        \langle u,v+w\rangle=\langle u,v\rangle+\langle u,w\rangle
    \]
    and
    \[
        \langle u,\lambda v\rangle=\lambda\langle u,v\rangle.
    \]
\end{snippetproposition}

\begin{snippetproof}{real-inner-product-second-variable-linearity-proof}{real-inner-product-second-variable-linearity}{Linearity in the second variable over the real numbers}
    By symmetry and linearity in the first variable,
    \[
        \langle u,v+w\rangle
        =
        \langle v+w,u\rangle
        =
        \langle v,u\rangle+\langle w,u\rangle
        =
        \langle u,v\rangle+\langle u,w\rangle.
    \]
    Similarly,
    \[
        \langle u,\lambda v\rangle
        =
        \langle \lambda v,u\rangle
        =
        \lambda\langle v,u\rangle
        =
        \lambda\langle u,v\rangle.
    \]
\end{snippetproof}

\begin{snippetproposition}{complex-inner-product-second-variable-antilinearity}{Antilinearity in the second variable over the complex numbers}
    Let \(V\) be a \vectorspace over \(\complexnumbers\), and let
    \(\langle\cdot,\cdot\rangle\) be an \innerproduct on \(V\). Then, for
    every \(u,v,w\in V\) and every \(\lambda\in\complexnumbers\),
    \[
        \langle u,v+w\rangle=\langle u,v\rangle+\langle u,w\rangle
    \]
    and
    \[
        \langle u,\lambda v\rangle=\overline{\lambda}\langle u,v\rangle.
    \]
\end{snippetproposition}

\begin{snippetproof}{complex-inner-product-second-variable-antilinearity-proof}{complex-inner-product-second-variable-antilinearity}{Antilinearity in the second variable over the complex numbers}
    By Hermitian symmetry and linearity in the first variable,
    \[
        \langle u,v+w\rangle
        =
        \overline{\langle v+w,u\rangle}
        =
        \overline{\langle v,u\rangle+\langle w,u\rangle}
        =
        \langle u,v\rangle+\langle u,w\rangle.
    \]
    Also,
    \[
        \langle u,\lambda v\rangle
        =
        \overline{\langle \lambda v,u\rangle}
        =
        \overline{\lambda\langle v,u\rangle}
        =
        \overline{\lambda}\langle u,v\rangle.
    \]
\end{snippetproof}

\begin{snippetdefinition}{bilinear-map-definition}{Bilinear map}
    Let \(V,W,Z\) be \vectorspace[vector spaces] over a \field
    \(\mathbb{F}\). A \function
    \[
        B\colon V\cartesianprod W\fromto Z
    \]
    is \emph{bilinear} if, for every \(v,v'\in V\), every \(w,w'\in W\),
    and every \(\lambda,\mu\in\mathbb{F}\),
    \[
        B(\lambda v+\mu v',w)
        =
        \lambda B(v,w)+\mu B(v',w)
    \]
    and
    \[
        B(v,\lambda w+\mu w')
        =
        \lambda B(v,w)+\mu B(v,w').
    \]
\end{snippetdefinition}

\begin{snippet}{real-inner-products-are-bilinear}
    Let \(V\) be a \vectorspace over \(\realnumbers\). An \innerproduct on
    \(V\) is a \bilinearmap
    \[
        \langle\cdot,\cdot\rangle\colon V\cartesianprod V\fromto\realnumbers
    \]
    that is symmetric and positive definite.
\end{snippet}

\begin{snippetdefinition}{inner-product-space-definition}{Inner product space}
    An \emph{inner product space} is a pair
    \[
        (V,\langle\cdot,\cdot\rangle),
    \]
    where \(V\) is a \vectorspace over a \field
    \(\mathbb{F}\in\{\realnumbers,\complexnumbers\}\), and
    \(\langle\cdot,\cdot\rangle\) is an \innerproduct on \(V\).
\end{snippetdefinition}

\section{Examples}

\begin{snippetexample}{canonical-inner-product-fn-example}{Canonical inner product on \(F^n\)}
    Let \(\mathbb{F}\in\{\realnumbers,\complexnumbers\}\). The canonical
    \innerproduct on \(\mathbb{F}^n\) is defined by
    \[
        \langle x,y\rangle
        =
        \sum_{i=1}^{n}x_i\overline{y_i},
    \]
    where
    \[
        x=(x_1,\ldots,x_n),
        \qquad
        y=(y_1,\ldots,y_n).
    \]
    If \(\mathbb{F}=\realnumbers\), this becomes
    \[
        \langle x,y\rangle=x_1y_1+\cdots+x_ny_n.
    \]
\end{snippetexample}

\begin{snippetexample}{matrix-trace-inner-product-example}{Trace inner product on matrices}
    Let \(m,n\in\naturalnumbers\). On the real \vectorspace
    \(\matrices_{m\times n}(\realnumbers)\), define
    \[
        \langle A,B\rangle=\trace(AB^\transpose).
    \]
    Then \(\langle\cdot,\cdot\rangle\) is an \innerproduct. Bilinearity
    follows from the linearity of matrix multiplication and the trace, and
    symmetry follows from
    \[
        \trace(AB^\transpose)=\trace(BA^\transpose).
    \]
    If \(A=(a_{ij})\), then
    \[
        \langle A,A\rangle
        =
        \trace(AA^\transpose)
        =
        \sum_{i=1}^{m}\sum_{j=1}^{n}a_{ij}^2.
    \]
    Hence \(\langle A,A\rangle\geq0\), and
    \(\langle A,A\rangle=0\) holds if and only if every entry of \(A\) is
    zero.
\end{snippetexample}

\begin{snippetexample}{polynomial-values-inner-product-example}{Values inner product on cubic polynomials}
    Let
    \[
        \mathcal{P}_3(\realnumbers)
        =
        \{p\in\realnumbers[x]\suchthat \polynomialdeg(p)\leq3\}\union\{0\}.
    \]
    For \(p,q\in\mathcal{P}_3(\realnumbers)\), define
    \[
        \langle p,q\rangle
        =
        p(0)q(0)+p(1)q(1)+p(2)q(2)+p(-1)q(-1).
    \]
    This is an \innerproduct on \(\mathcal{P}_3(\realnumbers)\). Indeed,
    it is bilinear and symmetric, and
    \[
        \langle p,p\rangle
        =
        p(0)^2+p(1)^2+p(2)^2+p(-1)^2\geq0.
    \]
    If \(\langle p,p\rangle=0\), then \(p\) has the four distinct roots
    \(0,1,2,-1\). Since \(p\) has degree at most \(3\), it follows that
    \(p=0\).
\end{snippetexample}

\begin{snippetexample}{polynomial-integral-inner-product-example}{Integral inner product on polynomials}
    On the real \vectorspace \(\realnumbers[x]\), define
    \[
        \langle p,q\rangle
        =
        \integral[0][1][p(t)q(t)][t].
    \]
    This is an \innerproduct. Bilinearity follows from the linearity of the
    integral, and symmetry is immediate. Moreover,
    \[
        \langle p,p\rangle
        =
        \integral[0][1][p(t)^2][t]\geq0.
    \]
    If \(\langle p,p\rangle=0\), then the continuous function \(p(t)^2\)
    is identically zero on \([0,1]\). Hence \(p\) has infinitely many roots,
    so \(p=0\).
\end{snippetexample}

\begin{snippetexample}{ell-two-inner-product-example}{Inner product on square-summable sequences}
    Let
    \[
        \ell^2(\realnumbers)
        =
        \left\{
            x=(x_n)_{n\geq1}
            \suchthat
            \sum_{n=1}^{+\infty}x_n^2<+\infty
        \right\}.
    \]
    For \(x,y\in\ell^2(\realnumbers)\), define
    \[
        \langle x,y\rangle
        =
        \sum_{n=1}^{+\infty}x_ny_n.
    \]
    The series is absolutely convergent because
    \[
        2|x_ny_n|\leq x_n^2+y_n^2
    \]
    for every \(n\). With this operation, \(\ell^2(\realnumbers)\) is an
    \innerproductspace.
\end{snippetexample}

\section{Induced norm}

\begin{snippetdefinition}{inner-product-induced-norm-definition}{Norm induced by an inner product}
    Let \((V,\langle\cdot,\cdot\rangle)\) be an \innerproductspace. The
    \emph{norm induced by the inner product} is the \function
    \[
        \|\cdot\|\colon V\fromto\realnumbers
    \]
    defined by
    \[
        \|v\|=\sqrt{\langle v,v\rangle}.
    \]
\end{snippetdefinition}

\begin{snippetproposition}{inner-product-induced-norm-basic-properties}{Basic properties of the induced norm}
    Let \((V,\langle\cdot,\cdot\rangle)\) be an \innerproductspace over
    \(\mathbb{F}\). Then, for every \(v\in V\) and every
    \(\lambda\in\mathbb{F}\),
    \begin{enumerate}
        \item \(\|v\|\geq0\);
        \item \(\|v\|=0\) \ifandonlyif \(v=0\);
        \item \(\|\lambda v\|=|\lambda|\|v\|\).
    \end{enumerate}
\end{snippetproposition}

\begin{snippetproof}{inner-product-induced-norm-basic-properties-proof}{inner-product-induced-norm-basic-properties}{Basic properties of the induced norm}
    The first two statements follow directly from positivity and
    definiteness of the \innerproduct. For homogeneity,
    \[
        \|\lambda v\|^2
        =
        \langle \lambda v,\lambda v\rangle
        =
        \lambda\overline{\lambda}\langle v,v\rangle
        =
        |\lambda|^2\|v\|^2.
    \]
    Taking square roots gives
    \[
        \|\lambda v\|=|\lambda|\|v\|.
    \]
\end{snippetproof}

\section{Orthogonality}

\begin{snippetdefinition}{orthogonal-vectors-definition}{Orthogonal vectors}
    Let \((V,\langle\cdot,\cdot\rangle)\) be an \innerproductspace. Two
    \vector[vectors] \(u,v\in V\) are \emph{orthogonal}, written
    \[
        u\perp v,
    \]
    if
    \[
        \langle u,v\rangle=0.
    \]
\end{snippetdefinition}

\begin{snippetproposition}{orthogonality-is-symmetric}{Orthogonality is symmetric}
    Let \((V,\langle\cdot,\cdot\rangle)\) be an \innerproductspace, and let
    \(u,v\in V\). Then
    \[
        u\perp v
        \qquad
        \Longleftrightarrow
        \qquad
        v\perp u.
    \]
\end{snippetproposition}

\begin{snippetproof}{orthogonality-is-symmetric-proof}{orthogonality-is-symmetric}{Orthogonality is symmetric}
    By Hermitian symmetry,
    \[
        \langle v,u\rangle=\overline{\langle u,v\rangle}.
    \]
    Hence \(\langle u,v\rangle=0\) if and only if
    \(\langle v,u\rangle=0\).
\end{snippetproof}

\begin{snippetproposition}{orthogonal-to-all-is-zero}{The only vector orthogonal to every vector}
    Let \((V,\langle\cdot,\cdot\rangle)\) be an \innerproductspace. If
    \(v\in V\) is \orthogonal to every \vector of \(V\), then
    \[
        v=0.
    \]
\end{snippetproposition}

\begin{snippetproof}{orthogonal-to-all-is-zero-proof}{orthogonal-to-all-is-zero}{The only vector orthogonal to every vector}
    Since \(v\) is \orthogonal to every \vector of \(V\), it is
    \orthogonal to itself. Thus
    \[
        \langle v,v\rangle=0.
    \]
    By definiteness of the \innerproduct, \(v=0\).
\end{snippetproof}

\begin{snippettheorem}{pythagorean-theorem-inner-product}{Pythagorean theorem}
    Let \((V,\langle\cdot,\cdot\rangle)\) be an \innerproductspace, and let
    \(u,v\in V\). If \(u\perp v\), then
    \[
        \|u+v\|^2=\|u\|^2+\|v\|^2.
    \]
\end{snippettheorem}

\begin{snippetproof}{pythagorean-theorem-inner-product-proof}{pythagorean-theorem-inner-product}{Pythagorean theorem}
    Since \(u\perp v\), both \(\langle u,v\rangle\) and
    \(\langle v,u\rangle\) are zero. Therefore
    \[
        \|u+v\|^2
        =
        \langle u+v,u+v\rangle
        =
        \langle u,u\rangle+\langle u,v\rangle+\langle v,u\rangle+\langle v,v\rangle
        =
        \|u\|^2+\|v\|^2.
    \]
\end{snippetproof}

\section{Projection and Cauchy-Schwarz}

\begin{snippetdefinition}{orthogonal-projection-definition}{Orthogonal projection onto a line}
    Let \((V,\langle\cdot,\cdot\rangle)\) be an \innerproductspace, let
    \(u,v\in V\), and suppose that \(v\neq0\). The \emph{orthogonal
    projection of \(u\) onto the line generated by \(v\)} is
    \[
        \operatorname{proj}_v(u)
        =
        \frac{\langle u,v\rangle}{\|v\|^2}v.
    \]
\end{snippetdefinition}

\begin{snippetproposition}{orthogonal-projection-decomposition}{Orthogonal projection decomposition}
    Let \((V,\langle\cdot,\cdot\rangle)\) be an \innerproductspace, let
    \(u,v\in V\), and suppose that \(v\neq0\). Then
    \[
        u=
        \operatorname{proj}_v(u)
        +
        \left(u-\operatorname{proj}_v(u)\right),
    \]
    where
    \[
        u-\operatorname{proj}_v(u)\perp v.
    \]
\end{snippetproposition}

\begin{snippetproof}{orthogonal-projection-decomposition-proof}{orthogonal-projection-decomposition}{Orthogonal projection decomposition}
    Put
    \[
        w=u-\frac{\langle u,v\rangle}{\|v\|^2}v.
    \]
    Then
    \[
        \langle w,v\rangle
        =
        \langle u,v\rangle
        -
        \frac{\langle u,v\rangle}{\|v\|^2}\langle v,v\rangle
        =
        \langle u,v\rangle-\langle u,v\rangle
        =
        0.
    \]
    Thus \(w\perp v\), and the displayed decomposition follows.
\end{snippetproof}

\begin{snippettheorem}{cauchy-schwarz-inequality}{Cauchy-Schwarz inequality}
    Let \((V,\langle\cdot,\cdot\rangle)\) be an \innerproductspace. Then,
    for every \(u,v\in V\),
    \[
        |\langle u,v\rangle|\leq\|u\|\|v\|.
    \]
    If \(u\neq0\) and \(v\neq0\), equality holds if and only if \(u\) and
    \(v\) are linearly dependent.
\end{snippettheorem}

\begin{snippetproof}{cauchy-schwarz-inequality-proof}{cauchy-schwarz-inequality}{Cauchy-Schwarz inequality}
    If \(v=0\), the inequality is immediate. Suppose that \(v\neq0\). By the
    \orthogonalprojection decomposition,
    \[
        u=
        \frac{\langle u,v\rangle}{\|v\|^2}v+w,
        \qquad
        w\perp v.
    \]
    The Pythagorean theorem gives
    \[
        \|u\|^2
        =
        \left\|
        \frac{\langle u,v\rangle}{\|v\|^2}v
        \right\|^2
        +
        \|w\|^2
        \geq
        \left\|
        \frac{\langle u,v\rangle}{\|v\|^2}v
        \right\|^2.
    \]
    Therefore
    \[
        \|u\|^2
        \geq
        \frac{|\langle u,v\rangle|^2}{\|v\|^4}\|v\|^2
        =
        \frac{|\langle u,v\rangle|^2}{\|v\|^2}.
    \]
    Multiplying by \(\|v\|^2\) and taking square roots yields
    \[
        |\langle u,v\rangle|\leq\|u\|\|v\|.
    \]

    If \(u,v\neq0\), equality holds if and only if \(\|w\|=0\), which is
    equivalent to \(w=0\). This means that \(u\) is a scalar multiple of
    \(v\). Conversely, if \(u\) and \(v\) are linearly dependent, the
    equality follows directly from homogeneity of the induced \norm.
\end{snippetproof}

\begin{snippettheorem}{inner-product-triangle-inequality}{Triangle inequality}
    Let \((V,\langle\cdot,\cdot\rangle)\) be an \innerproductspace. Then,
    for every \(u,v\in V\),
    \[
        \|u+v\|\leq\|u\|+\|v\|.
    \]
\end{snippettheorem}

\begin{snippetproof}{inner-product-triangle-inequality-proof}{inner-product-triangle-inequality}{Triangle inequality}
    We have
    \[
        \|u+v\|^2
        =
        \|u\|^2+\langle u,v\rangle+\langle v,u\rangle+\|v\|^2.
    \]
    Since
    \[
        \langle u,v\rangle+\langle v,u\rangle
        =
        2\operatorname{Re}\langle u,v\rangle
        \leq
        2|\langle u,v\rangle|,
    \]
    the \cauchyschwarz gives
    \[
        \|u+v\|^2
        \leq
        \|u\|^2+2\|u\|\|v\|+\|v\|^2
        =
        (\|u\|+\|v\|)^2.
    \]
    Taking square roots gives the result.
\end{snippetproof}

\section{Induced metric}

\begin{snippetdefinition}{inner-product-induced-metric-definition}{Metric induced by an inner product}
    Let \((V,\langle\cdot,\cdot\rangle)\) be an \innerproductspace. The
    \emph{metric induced by the inner product} is the \function
    \[
        d\colon V\cartesianprod V\fromto\realnumbers
    \]
    defined by
    \[
        d(u,v)=\|u-v\|.
    \]
\end{snippetdefinition}

\begin{snippetproposition}{inner-product-induced-metric-is-metric}{The induced metric is a metric}
    Let \((V,\langle\cdot,\cdot\rangle)\) be an \innerproductspace. The
    \inducedmetric \(d(u,v)=\|u-v\|\) satisfies, for every \(u,v,w\in V\),
    \begin{enumerate}
        \item \(d(u,v)\geq0\);
        \item \(d(u,v)=0\) \ifandonlyif \(u=v\);
        \item \(d(u,v)=d(v,u)\);
        \item \(d(u,w)\leq d(u,v)+d(v,w)\).
    \end{enumerate}
\end{snippetproposition}

\begin{snippetproof}{inner-product-induced-metric-is-metric-proof}{inner-product-induced-metric-is-metric}{The induced metric is a metric}
    Positivity follows from positivity of the induced \norm. Moreover,
    \[
        d(u,v)=0
        \Longleftrightarrow
        \|u-v\|=0
        \Longleftrightarrow
        u-v=0
        \Longleftrightarrow
        u=v.
    \]
    Symmetry follows from homogeneity:
    \[
        d(u,v)=\|u-v\|=\|-(v-u)\|=\|v-u\|=d(v,u).
    \]
    Finally, the triangle inequality for the induced \norm gives
    \[
        d(u,w)
        =
        \|u-w\|
        =
        \|(u-v)+(v-w)\|
        \leq
        \|u-v\|+\|v-w\|
        =
        d(u,v)+d(v,w).
    \]
\end{snippetproof}

\section{Angles}

\begin{snippetdefinition}{angle-between-vectors-definition}{Angle between two vectors}
    Let \((V,\langle\cdot,\cdot\rangle)\) be a real \innerproductspace, and
    let \(u,v\in V\) be nonzero \vector[vectors]. The \emph{angle between
    \(u\) and \(v\)} is the unique number \(\theta\in[0,\pi]\) such that
    \[
        \cos\theta
        =
        \frac{\langle u,v\rangle}{\|u\|\|v\|}.
    \]
\end{snippetdefinition}

\begin{snippetproposition}{angle-between-vectors-well-defined}{The angle between vectors is well-defined}
    Let \((V,\langle\cdot,\cdot\rangle)\) be a real \innerproductspace, and
    let \(u,v\in V\) be nonzero \vector[vectors]. Then there exists a unique
    \(\theta\in[0,\pi]\) such that
    \[
        \cos\theta
        =
        \frac{\langle u,v\rangle}{\|u\|\|v\|}.
    \]
\end{snippetproposition}

\begin{snippetproof}{angle-between-vectors-well-defined-proof}{angle-between-vectors-well-defined}{The angle between vectors is well-defined}
    By the \cauchyschwarz,
    \[
        |\langle u,v\rangle|\leq\|u\|\|v\|.
    \]
    Since \(u\) and \(v\) are nonzero, \(\|u\|\|v\|>0\), and therefore
    \[
        -1
        \leq
        \frac{\langle u,v\rangle}{\|u\|\|v\|}
        \leq
        1.
    \]
    The function
    \[
        \cos\colon[0,\pi]\fromto[-1,1]
    \]
    is bijective, so the required \(\theta\) exists and is unique.
\end{snippetproof}

\begin{snippetproposition}{orthogonality-angle-pi-over-two}{Orthogonality and right angles}
    Let \((V,\langle\cdot,\cdot\rangle)\) be a real \innerproductspace, and
    let \(u,v\in V\) be nonzero \vector[vectors]. Then
    \[
        u\perp v
        \qquad
        \Longleftrightarrow
        \qquad
        \theta(u,v)=\frac{\pi}{2},
    \]
    where \(\theta(u,v)\) is the angle between \(u\) and \(v\).
\end{snippetproposition}

\begin{snippetproof}{orthogonality-angle-pi-over-two-proof}{orthogonality-angle-pi-over-two}{Orthogonality and right angles}
    By definition of the angle,
    \[
        \cos\theta(u,v)
        =
        \frac{\langle u,v\rangle}{\|u\|\|v\|}.
    \]
    Since \(u\) and \(v\) are nonzero, the denominator is positive. Thus
    \(\langle u,v\rangle=0\) if and only if
    \(\cos\theta(u,v)=0\). On \([0,\pi]\), this happens if and only if
    \[
        \theta(u,v)=\frac{\pi}{2}.
    \]
\end{snippetproof}

\section{Associated matrix}

\begin{snippetdefinition}{inner-product-associated-matrix-definition}{Matrix associated to an inner product}
    Let \((V,\langle\cdot,\cdot\rangle)\) be a finite-dimensional
    \innerproductspace over \(\mathbb{F}\in\{\realnumbers,\complexnumbers\}\),
    with \(\lineardim V=n\). Let
    \[
        \mathcal{B}=\{v_1,\ldots,v_n\}
    \]
    be a \basis of \(V\). The \emph{matrix associated to the inner product
    in the basis \(\mathcal{B}\)} is the \matrix
    \[
        S_{\mathcal{B}}=(s_{ij})\in\matrices_{n\times n}(\mathbb{F})
    \]
    defined by
    \[
        s_{ij}=\langle v_i,v_j\rangle.
    \]
\end{snippetdefinition}

\begin{snippetproposition}{inner-product-coordinate-formula}{Coordinate formula for an inner product}
    Let \((V,\langle\cdot,\cdot\rangle)\) be a finite-dimensional
    \innerproductspace over \(\mathbb{F}\in\{\realnumbers,\complexnumbers\}\),
    let \(\mathcal{B}\) be a \basis of \(V\), and let
    \(S_{\mathcal{B}}\) be the \innerproductmatrix of
    \(\langle\cdot,\cdot\rangle\) in \(\mathcal{B}\). If \(u,w\in V\), then
    \[
        \langle u,w\rangle
        =
        [u]_{\mathcal{B}}^\transpose
        S_{\mathcal{B}}
        \overline{[w]_{\mathcal{B}}}.
    \]
    In the real case this becomes
    \[
        \langle u,w\rangle
        =
        [u]_{\mathcal{B}}^\transpose
        S_{\mathcal{B}}
        [w]_{\mathcal{B}}.
    \]
\end{snippetproposition}

\begin{snippetproof}{inner-product-coordinate-formula-proof}{inner-product-coordinate-formula}{Coordinate formula for an inner product}
    Write
    \[
        \mathcal{B}=\{v_1,\ldots,v_n\},
        \qquad
        u=\sum_{i=1}^{n}\alpha_i v_i,
        \qquad
        w=\sum_{j=1}^{n}\beta_j v_j.
    \]
    By linearity in the first variable and conjugate-linearity in the second
    variable,
    \[
        \langle u,w\rangle
        =
        \sum_{i=1}^{n}\sum_{j=1}^{n}
        \alpha_i\overline{\beta_j}\langle v_i,v_j\rangle.
    \]
    This is exactly
    \[
        [u]_{\mathcal{B}}^\transpose
        S_{\mathcal{B}}
        \overline{[w]_{\mathcal{B}}}.
    \]
    If \(\mathbb{F}=\realnumbers\), conjugation is trivial.
\end{snippetproof}

\section{Symmetric, Hermitian and positive definite matrices}

\includesnpt{symmetric-skew-symmetric-matrix-definition}

\includesnpt{hermitian-matrix-definition}

\begin{snippetdefinition}{positive-definite-matrix-definition}{Positive definite matrix}
    Let \(\mathbb{F}\in\{\realnumbers,\complexnumbers\}\), and let
    \(S\in\matrices_{n\times n}(\mathbb{F})\). The \matrix \(S\) is
    \emph{positive definite} if, for every nonzero
    \(x\in\mathbb{F}^n\),
    \[
        x^\transpose S\overline{x}>0.
    \]
    In the real case this condition is
    \[
        x^\transpose Sx>0
    \]
    for every nonzero \(x\in\realnumbers^n\).
\end{snippetdefinition}

\begin{snippetproposition}{inner-product-matrix-properties}{Properties of the matrix associated to an inner product}
    Let \((V,\langle\cdot,\cdot\rangle)\) be a finite-dimensional
    \innerproductspace over \(\mathbb{F}\in\{\realnumbers,\complexnumbers\}\),
    and let \(\mathcal{B}\) be a \basis of \(V\). The \innerproductmatrix
    \(S_{\mathcal{B}}\) is:
    \begin{enumerate}
        \item a \symmetricmatrix if \(\mathbb{F}=\realnumbers\);
        \item a \hermitianmatrix if \(\mathbb{F}=\complexnumbers\);
        \item a \positivedefinitematrix.
    \end{enumerate}
\end{snippetproposition}

\begin{snippetproof}{inner-product-matrix-properties-proof}{inner-product-matrix-properties}{Properties of the matrix associated to an inner product}
    Let
    \[
        \mathcal{B}=\{v_1,\ldots,v_n\},
        \qquad
        S_{\mathcal{B}}=(s_{ij}).
    \]
    If \(\mathbb{F}=\realnumbers\), then
    \[
        s_{ij}
        =
        \langle v_i,v_j\rangle
        =
        \langle v_j,v_i\rangle
        =
        s_{ji},
    \]
    so \(S_{\mathcal{B}}\) is a \symmetricmatrix. If
    \(\mathbb{F}=\complexnumbers\), then
    \[
        s_{ij}
        =
        \langle v_i,v_j\rangle
        =
        \overline{\langle v_j,v_i\rangle}
        =
        \overline{s_{ji}},
    \]
    so \(S_{\mathcal{B}}\) is a \hermitianmatrix.

    Finally, let \(x\in\mathbb{F}^n\) be nonzero, and let \(v\in V\) be the
    unique \vector such that \([v]_{\mathcal{B}}=x\). Since \(x\neq0\), we
    have \(v\neq0\). By the coordinate formula,
    \[
        x^\transpose S_{\mathcal{B}}\overline{x}
        =
        \langle v,v\rangle
        >
        0.
    \]
    Hence \(S_{\mathcal{B}}\) is a \positivedefinitematrix.
\end{snippetproof}

\section{Change of basis}

\begin{snippetproposition}{inner-product-matrix-change-of-basis}{Change of basis for the inner product matrix}
    Let \((V,\langle\cdot,\cdot\rangle)\) be a finite-dimensional
    \innerproductspace over \(\mathbb{F}\in\{\realnumbers,\complexnumbers\}\).
    Let \(\mathcal{B}\) and \(\mathcal{C}\) be \basis[bases] of \(V\), and
    let \(S_{\mathcal{B}}\) and \(S_{\mathcal{C}}\) be the
    \innerproductmatrix[associated matrices] of the same \innerproduct in
    the two bases. If
    \[
        A=M(\identityfunc_V,\mathcal{B},\mathcal{C}),
    \]
    then
    \[
        S_{\mathcal{B}}
        =
        A^\transpose S_{\mathcal{C}}\overline{A}.
    \]
    In the real case this becomes
    \[
        S_{\mathcal{B}}
        =
        A^\transpose S_{\mathcal{C}}A.
    \]
\end{snippetproposition}

\begin{snippetproof}{inner-product-matrix-change-of-basis-proof}{inner-product-matrix-change-of-basis}{Change of basis for the inner product matrix}
    By definition of the change-of-basis \matrix,
    \[
        [v]_{\mathcal{C}}=A[v]_{\mathcal{B}}
    \]
    for every \(v\in V\). Therefore, for every \(u,v\in V\),
    \[
        \langle u,v\rangle
        =
        [u]_{\mathcal{B}}^\transpose
        S_{\mathcal{B}}
        \overline{[v]_{\mathcal{B}}}
    \]
    and also
    \[
        \langle u,v\rangle
        =
        [u]_{\mathcal{C}}^\transpose
        S_{\mathcal{C}}
        \overline{[v]_{\mathcal{C}}}
        =
        [u]_{\mathcal{B}}^\transpose
        A^\transpose S_{\mathcal{C}}\overline{A}
        \overline{[v]_{\mathcal{B}}}.
    \]
    Since this holds for every pair of coordinate vectors, the matrices are
    equal:
    \[
        S_{\mathcal{B}}
        =
        A^\transpose S_{\mathcal{C}}\overline{A}.
    \]
    If \(\mathbb{F}=\realnumbers\), conjugation is trivial.
\end{snippetproof}

\begin{snippetdefinition}{congruent-matrices-definition}{Congruent matrices}
    Let \(\mathbb{F}\in\{\realnumbers,\complexnumbers\}\). Two
    \matrix[matrices]
    \[
        S,T\in\matrices_{n\times n}(\mathbb{F})
    \]
    are \emph{congruent} if there exists an invertible
    \[
        A\in\matrices_{n\times n}(\mathbb{F})
    \]
    such that
    \[
        S=A^\transpose T\overline{A}.
    \]
    If \(\mathbb{F}=\realnumbers\), this condition becomes
    \[
        S=A^\transpose TA.
    \]
\end{snippetdefinition}

\begin{snippet}{real-inner-product-matrices-are-congruent-under-basis-change}
    Let \(V\) be a finite-dimensional real \innerproductspace. If
    \(S_{\mathcal{B}}\) and \(S_{\mathcal{C}}\) are the
    \innerproductmatrix[associated matrices] of the same \innerproduct in two
    \basis[bases] \(\mathcal{B}\) and \(\mathcal{C}\), then
    \(S_{\mathcal{B}}\) and \(S_{\mathcal{C}}\) are \congruentmatrices.
\end{snippet}

\end{document}
